\documentclass[11pt]{article}
\usepackage{amsmath,amsthm,amssymb}
\usepackage[margin=1in]{geometry}
\usepackage{booktabs}
\newtheorem{theorem}{Theorem}
\newtheorem{proposition}[theorem]{Proposition}
\newtheorem{lemma}[theorem]{Lemma}
\theoremstyle{definition}
\newtheorem{definition}[theorem]{Definition}
\newtheorem{remark}[theorem]{Remark}
\title{Q585 re-aimed: Purbhoo's flock swap equals the crystal commutor\\
(computed on 190 cells, $c^\nu_{\lambda\mu}\in\{2,3,4,5\}$, under two bijections)}
\author{Clio}
\date{2026-11-12, PROVE session c2.\\[2pt]\small To: Robin. From: Clio.
Provenance (PROTOCOL \S2.3): sources commit \texttt{7303d1db5f52} (clio-vega/proofs, \texttt{code-1112c2-Q585/}); addendum \S\ref{sec:split}: \texttt{3257a1d79e78} (\texttt{switch.py});
C instrument from \texttt{code-1112-Q585/} (PROVE c1). Source read first-hand: Purbhoo, arXiv:0705.1184, \S\S2--4 and \S5.1
(text, and the figures on pp.~6 and~13 rendered at 400--900\,dpi). \textbf{Grade: computed, not proved.}}
\begin{document}
\maketitle

\begin{abstract}
Purbhoo proves commutativity of the Littlewood--Richardson rule bijectively by migrating flocks of
$(150^\circ,30^\circ)$ rhombi through a mosaic \cite[Cor.~3.4]{P}. The result is a map
$M:\mathrm{LR}(\lambda/\mu;\nu)\to\mathrm{LR}(\lambda/\nu;\mu)$, which he states is an involution
\cite[Prop.~5.4]{P}. I built migration from the definitions, exactly and from scratch. I then compared $M$ with the
Henriques--Kamnitzer crystal commutor $C$ on every cell with $c\ge2$ in the range
$\mathfrak{gl}_4$, $|\lambda|,|\mu|\le7$ ($c\ge3$) or $\le5$ ($c=2$). \textbf{They agree on $518/518$ elements in
$190/190$ cells, under each of two independent tableau$\leftrightarrow$highest-weight bijections.} The lex control
separates from $C$ in every cell. This is computed evidence and not a proof. The gap is stated in \S6.
\end{abstract}

\section{The question and the pre-registration}
The brief (WAKE 2026-11-12 c2) asked whether Purbhoo's swap equals $C$ on highest-weight elements. It pre-registered
``I don't know; both outcomes informative'' and named the falsifier as one $c\ge2$ cell where the swap differs from
$C$, under a bijection that passes its own round-trip control. The falsifier did not fire.

\section{Definitions, as read (\cite{P} \S\S2.3--3.3)}
\begin{enumerate}
\item \emph{Mosaic.} A mosaic is a tiling of the hexagon $A'AB'BC'C$ (angles $150^\circ$ at $A,B,C$ and $90^\circ$ at $A',B',C'$)
by unit triangles, unit squares and $(150^\circ,30^\circ)$ rhombi, with the rhombi packed into nests at $A,B,C$. Mosaics
are in bijection with KT puzzles: delete the rhombi and flex the squares to $60^\circ$ rhombi; the $0$-edges and
$1$-edges rotate in opposite senses. The standard orientations are $(E_A,N_A)=(\vec{AA'},\vec{AB'})$,
$(E_B,N_B)=(\vec{BB'},\vec{BC'})$, $(E_C,N_C)=(\vec{CC'},\vec{CA'})$.
\item \emph{Flock.} A flock is $(\zeta,(E,N),f)$, where $\zeta$ is a set of nest rhombi forming a transformed skew
diagram and $f$ is an LR tableau on $\zeta$ in the frame $(E,N)$. The \emph{standard order} compares $f$-values first
and breaks ties by ``west first''.
\item \emph{Migration} (\S3.1). Turn the page so that the direction of migration points right. N/E migration processes
the rhombi from the standard-order maximum downward, and S/W migration from the minimum upward. Each rhombus
$\lozenge$ repeatedly finds the smallest hexagon containing it that lies weakly right of its leftmost point.
Figures (a) and (b) on p.~13 show both kinds as \emph{rhombus $+$ square $+$ two triangles}. If the hexagon is
$2$-fold symmetric (a), rotate its tiling by $180^\circ$. If it is $3$-fold symmetric (b), rotate by $120^\circ$ in
the sense that leaves an edge of $\lozenge$ horizontal or vertical. Stop when $\lozenge$ is packed in a new nest.
\item \emph{The swap} (Cor.~3.4). Let $f$ be an LR tableau on $\lambda/\mu$ with content $\nu$. Form the canonical mosaic
$(\varnothing,\gamma^\vee,\gamma)$ with $\gamma^\vee=\lambda$ in $(E_B,N_B)$ (Lemma 2.6). Let $F$ be the flock of $f$,
and $G$ the unique flock on the remaining rhombi, in orientation $(-E_B,-N_B)$. Migrate $F$ from $B$ to $A$, then
$G$ from $B$ to $A$. Then $G'$ is an LR tableau on $\lambda/\nu$ with content $\mu$.
\end{enumerate}
\paragraph{What the paper does not make constructive.} Lemma~2.6 (the canonical mosaic) is stated with ``proofs
omitted''; I build it from the unique KT puzzle with that boundary and assert uniqueness. The proofs of Prop.~5.4 and
of the jdt dictionary (Props.~5.1--5.2) are omitted. Cor.~3.3 is a \emph{family} of swaps (``fix any orientations''),
so ``the'' Purbhoo swap is well-defined only for Cor.~3.4, which is what is tested here. Text-level readings:
Cor.~3.4 says $\eta$ ``corresponds to $\nu$'', but the remaining rhombi are $\mu$, so I read it as $\mu$;
\S2.2(i) says ``east to west'', but Fig.~9 increases west to east, and I follow Fig.~9.

\section{The instrument}
\paragraph{Exact geometry.} A point is the KT lift $P=(P_0,P_1)\in\mathbb Z[\omega]^2\cong\mathbb Z^4$, the sums of the
$0$-labelled and $1$-labelled edge vectors along any path, as in Knutson--Tao's $\mathbb R^4$ surface (\cite[\S2.4]{P}).
It is drawn as $\phi(P)=e^{i\pi/12}P_0+e^{-i\pi/12}P_1$. The map $\phi$ is injective, since $e^{i\pi/6}\notin\mathbb Q(\omega)$.
Rotations by $60^\circ k$ act componentwise and exactly, so every tile coordinate is an integer vector. Floats are used
only for the left/right and horizontal tests. Of the four sign/corner conventions, \emph{exactly one} makes every KT
rhombus a square and every mosaic convex and area-exact ($41/41$ at $n=4$). So the convention is forced, not fitted.

\paragraph{Corrections found by the instrument, in order.} (i) My first low-resolution reading of Fig.~(a) took it as
three tiles. At 900\,dpi it is four: both moves are rhombus $+$ square $+$ two triangles. (ii) The union's boundary has a
straight-angle vertex ($30+90+60=180$), so ``hexagon'' means six \emph{corners} on seven unit edges. (iii) In 956
steps two hexagons tie for minimal. Both lie weakly right of the leftmost point, which contradicts the paper's uniqueness
parenthesis. I break the tie by the pattern visible in both figures: \emph{the hexagon contains the tile across the
horizontal edge of $\lozenge$ that starts at its leftmost vertex.} This is an interpretation, and it is flagged. It
is load-bearing: the lesion that takes the other candidate gets stuck on $63/142$ mosaics at $n=5$. With the rule,
all ties resolve ($956/956$).

\paragraph{Validation against the paper's own theorems}, before Cor.~3.4 was touched. On all mosaics with $n=4,5$
($27+112$), for all $8$ orientation/direction combinations of Table~1, the unique flock on $\alpha$ migrates
\emph{every} time. It lands in the nest the text predicts (for example $(-E_A,-N_A)$ south $\to B$, as in Cor.~3.2),
and the result is an LR tableau (Thm.~1) in $100\%$ of cases. Cor.~3.2 is injective within each boundary: $142$
mosaics, $0$ collisions.

\section{Transport: two bijections}
$C$ acts on $\mathrm{hw}(B(\lambda)\otimes B(\mu))_\nu$, realised as $b\in\mathrm{SSYT}(\lambda)$ with
$\mathrm{read}(b)\,\mathrm{read}(u_\mu)$ Yamanouchi. It is computed by two instruments from c1 (crystal operators;
RSK plus evacuation), which agree on $1308/1308$ elements.
\begin{itemize}
\item $\beta_{\rm CM}$: the content-matrix transpose. Row $r$ of $\nu/\mu$ holds $m_{k,r}$ copies of $k$, where $m_{k,r}=\#\{r{+}1\text{ in row }k\text{ of }b\}$.
This gives $\mathrm{hw}\to\mathrm{LR}(\nu/\mu;\lambda)$.
\item $\beta_Q$: the recording tableau. Write $w=\mathrm{read}(b)\,w'$, where $w'=\mathrm{RSK}^{-1}(u_\mu,U_\mu)$ and $U_\mu$
is row-superstandard, so that $w'\equiv u_\mu$ (Knuth) with a monotone destandardisation. The entries of $Q(w)$ above
$|\lambda|$ form a standard skew tableau on $\nu/\lambda$ that rectifies to $U_\mu$. Relabel each entry by its row in
$U_\mu$. This gives $\mathrm{hw}\to\mathrm{LR}(\nu/\lambda;\mu)$. It is valid and injective on $1308/1308$.
Two na\"ive versions failed their own round-trip control and were discarded: recording the inserted letter gives
anti-LR tableaux ($990/1308$ bad), and a row-reading relabel is non-monotone ($326/1308$ bad).
\end{itemize}
The two bijections land in \emph{different} sets: $\beta_{\rm CM}$ feeds $M$ an input on $\nu/\mu$, and $\beta_Q$ one
on $\nu/\lambda$. So they test $M$ on different inputs.

\section{Result}
\begin{proposition}[computed]\label{prop:main}
For every $\mathfrak{gl}_4$ cell with $c^\nu_{\lambda\mu}\ge2$ and $|\lambda|,|\mu|\le5$, and every cell with
$c^\nu_{\lambda\mu}\ge3$ and $|\lambda|,|\mu|\le7$, and for $\beta\in\{\beta_Q,\beta_{\rm CM}\}$,
\[ M\circ\beta \;=\; \beta\circ C \quad\text{on }\mathrm{hw}(B(\lambda)\otimes B(\mu))_\nu. \]
\end{proposition}
\begin{center}\begin{tabular}{rrrrrr}\toprule
$c$ & cells & elements & agree ($\beta_Q$) & agree ($\beta_{\rm CM}$) & lex $\ne C$ \\ \midrule
2 & 67 & 134 & 134 & 134 & 67/67\\
3 & 109 & 327 & 327 & 327 & 109/109\\
4 & 13 & 52 & 52 & 52 & 13/13\\
5 & 1 & 5 & 5 & 5 & 1/1\\ \midrule
all & 190 & 518 & 518 & 518 & 190/190\\ \bottomrule
\end{tabular}\end{center}
In every case the output of $M$ is an LR tableau of the predicted shape (Thm.~1 at scale), and $M\circ M=\mathrm{id}$
(Prop.~5.4 confirmed, $0$ failures under either bijection).

\paragraph{How many bits this is.} In a $c=2$ cell there are only two bijections, and the c1 controls showed that
lex $\neq C$ and reverse-lex $=C$ there. So on the $67$ $c=2$ cells, ``$M\ne$ lex'' and ``$M=C$'' are \emph{the same
bit} (the c1 lesson: a cheap control can be the target in disguise). Those cells carry $67$ bits, not more. The $c\ge3$
cells are the real test. A $c=3$ cell has $6$ bijections and a $c=4$ cell has $24$, and $M$ hits $C$'s in all $123$.
There, lex differs from $C$ in $123/123$ cells, but reverse-lex differs in only $27/123$. So reverse-lex remains a
\emph{near-blind} control, and the separation from it rests on those $27$ cells.

\paragraph{A lesion that must fire.} Replace $G$'s orientation $(-E_B,-N_B)$ by $(E_B,N_B)$ and migrate it east. This
is a legal migration but not Purbhoo's recipe. It differs from $C$ on $67/67$ $c=2$ cells. It is a crude lesion: $111$
of its outputs are not even LR in that frame.

\section{What is and is not established}
\begin{itemize}
\item \textbf{Computed:} Proposition~\ref{prop:main}, with the flagged tie-break rule as part of the definition of
migration. The verdict does not depend on which of the two bijections is used.
\item \textbf{Not proved.} A proof would need two steps. (1) A proof that Cor.~3.4's migrations are a specific jdt
composite. Purbhoo's \S5.1 asserts that migration corresponds to jdt (Props.~5.1, 5.2, Cor.~5.3), but omits the proofs,
and the B$\to$A migrations used in Cor.~3.4 are not covered by the stated B$\to$C case. (2) An identification of that
composite with $C=\xi_{B(\mu)\otimes B(\lambda)}\circ(\xi\otimes\xi)\circ\mathrm{flip}$. The natural candidate is a
tableau-switching description of $C$. I recall such an identification in the literature (Benkart--Sottile--Stroomer
switching, and an equivalence result for LR symmetry maps), but \emph{from memory only, not read}. So it is not cited
here and nothing rests on it.
\item \textbf{Consequence for Q583/PV05} (\emph{hunch-grade}, inherited framing): the ``migration train'' at the Schur
point induces the coboundary commutor on highest-weight elements. On this axis PV05 is not a rival \emph{involution}:
same map, different construction. That says nothing about the YBE/R-matrix train of c1, which remains ill-posed at the
instrument.
\end{itemize}

\section{Addendum (same session, re-entry): the gap splits through tableau switching}\label{sec:split}
\paragraph{The intermediate map.} For an LR tableau $S$ on $\nu/\kappa$ of content $\rho$, define $\Sigma(S)$ as follows.
Place $u_\kappa$ on $\kappa$. Take the cells of $\kappa$ in decreasing standard order of $u_\kappa$: larger entry
first, equal entries right to left. Slide $S$ inward into each cell by ordinary jeu de taquin (the hole takes the
smaller of its right and next-row neighbours; on a tie it takes the next-row one). Label the outer cell vacated
by that slide with the $u_\kappa$-entry of the starting cell. The labels form $\Sigma(S)$ on $\nu/\rho$. This is
switching of the pair $(u_\kappa,S)$ in the sense of Benkart--Sottile--Stroomer. The name is recalled, not read, and
nothing below rests on any property of switching from the literature: every property used is checked by computer.

\begin{proposition}[computed; same census as Prop.~\ref{prop:main}: 190 cells, 518 elements]\label{prop:split}
\begin{enumerate}
\item[(H1)] $M(T)=\Sigma(T)$ for every $T=\beta_{\rm CM}(b)\in\mathrm{LR}(\nu/\mu;\lambda)$: $518/518$.
\item[(H2)] $\Sigma\circ\beta=\beta\circ C$ for \emph{both} $\beta=\beta_{\rm CM}$ ($518/518$) and $\beta=\beta_Q$ ($518/518$).
\end{enumerate}
Moreover, $\Sigma(T)$ is always LR of content $\mu$ ($518/518$), and $\Sigma\circ\Sigma=\mathrm{id}$ ($518/518$).
\end{proposition}
\paragraph{Controls.} $\Sigma$ differs from lex in $190/190$ cells. Two lesions of $\Sigma$, namely breaking ties toward
the right neighbour and processing $\kappa$ column-first, together differ from $\Sigma$ in $190/190$ cells. So $\Sigma$
is not an order-insensitive artefact.

\paragraph{What this buys.} H1 and H2 together imply Prop.~\ref{prop:main}. Each is a separate claim, and so the single
open gap becomes two gaps of different kinds:
\begin{itemize}
\item \textbf{H1 is the Purbhoo-specific half.} It says that the two B$\to$A migrations of Cor.~3.4 together
\emph{are} switching against $u_\lambda$. A proof would go through Purbhoo's jdt dictionary (\S5.1). That dictionary
has its proofs omitted and is stated only for B$\to$C. So H1 is a \emph{proof} task about mosaics.
\item \textbf{H2 is Purbhoo-free.} It identifies switching with the HK commutor on highest-weight elements. I hold this
in the literature only at agent-summary grade (e.g.\ \texttt{0906.0077}), so H2 is a \emph{reading} task. It is not cited
here.
\end{itemize}
Neither half is proved in this note.

\section{Verification artifacts}
\texttt{proofs/code-1112c2-Q585/}: \texttt{mosaic.py} (puzzles, lift, mosaics), \texttt{migrate.py} (moves,
journeys), \texttt{flock.py} (flocks, frames, Cor.~3.4), \texttt{geotest.py} (convention 1 of 4),
\texttt{cor32.py}/\texttt{cor32b.py} (Table-1 validation, injectivity, tie-break lesion),
\texttt{betaq.py}, \texttt{compare.py}, \texttt{compare2.py} (the result), \texttt{ctl3.py} (controls),
\texttt{lesion34.py}, \texttt{switch.py} (Prop.~\ref{prop:split}), \texttt{results-*.txt}, \texttt{DEFS.md} (the definition summary written before coding).

\begin{thebibliography}{9}
\bibitem{P} K.~Purbhoo, \emph{Puzzles, tableaux, and mosaics}, arXiv:0705.1184 (J.~Algebraic Combin.\ 2008), \S\S2--5.
\bibitem{HK} A.~Henriques, J.~Kamnitzer, \emph{Crystals and coboundary categories}, arXiv:math/0406478 (the commutor
formula used, as implemented in c1; formula quoted via \texttt{0707.2248} l.855 in c1).
\end{thebibliography}
\end{document}
